\section{适用冲突规范的制度}
\begin{enumerate}
  \item 六个重要的制度，重在理解，\zd{外国法的查明、直接适用的法}（考试重复率最高的）。
  \item 定性（识别）。
        \begin{enumerate}
          \item 是对案件事实调查的过程，还是性质认定的过程。
          \item 分为大定性和小定性。
          \item \zd{法院地法识别说}。
          \item 识别的分割制：涉及两个 \hz 两个以上的法律关系\zd{分别适用}。
        \end{enumerate}
  \item 反致：一波三折，折腾。
        \begin{kk}{分类}
          \begin{enumerate}
            \item 转致：甲 \yjt 乙 \yjt 丙。折腾出去了。
            \item 直接反致：甲 \yjt 乙 \yjt 甲。折腾回来了。
            \item 间接反致：甲 \yjt 乙 \yjt 丙 \yjt 甲。折腾一圈才回来。
            \item 包含直接反致的转致：甲 \yjt 乙 \yjt 丙 \yjt 乙。回溯。
          \end{enumerate}
          \tcblower
          用了非法院地法作判决 \yjt 转致，法院地 \yjt 反致。

          \zd{我国司法实践中禁止反致与转致，选第一个国家}。
        \end{kk}
  \item 外国法查明和适用。
        \begin{enumerate}
          \item 当事人查 \yjt 谁审谁查 \yjt \zd{无法查明的} \yjt 中国法 \yjt
                \begin{kk}{无法查明}
                  \begin{itemize}
                    \item 当事人合理期限内，无正当理由，不能提供，视为无法查明。\zd{查不来}。
                    \item 穷尽所有途径：当事人、专家/机构、司法协助协定、我国驻该国使领馆、\zd{该国驻我国使馆}。
                  \end{itemize}
                  \tcblower
                  \qd{当家机协使领馆，外驻我领非途径}。
                \end{kk}
          \item 查来了 \yjt 当事人有异议 \yjt 听取当事人意见 \yjt 审查认定 \yjt 无误之后适用。
          \item \qd{不是证据的提供方，不是证人、鉴定人、专家}。
        \end{enumerate}
  \item 公共秩序保留：个案中认定，弹性比较大，考得不多，\zd{结果：排除外国法，适用中国法}。
  \item \zd{法律规避}。
        \begin{enumerate}
          \item 当事人\zd{故意}。
          \item \zd{制造 \hz 变更连结点}。
          \item 避开本应适用的\zd{对其不利}的准据法。
          \item 使\zd{对其有利}的法律得以适用。
          \item 第 c 点：我国只是不允许\zd{规避中国法} \hz \zd{法律 \hz 行政法规}的\zd{强制性规定}。热点，没考过。
        \end{enumerate}
  \item 直接适用中国的法：司法解释进行了限缩性解释，范围予以了限制，\zd{一保护两反三安全}。
        \begin{enumerate}
          \item 劳动者权益保护。\zd{题目会明确指明}。
          \item 反垄断反倾销。
          \item 食品 \hz 公共卫生安全、环境安全、外汇管制等金融安全。
        \end{enumerate}
\end{enumerate}
